\writeSectionFrame{\presentationTitle}{Fazit}
\begin{frame}{Was wir gelernt haben}
	\begin{itemize}
	    \item Dependency Injection entkoppelt Klassen von konkreten Implementierungen
	    \item Constructor Injection: Abhängigkeiten werden bei Erstellung gesetzt
	    \item Setter Injection: Abhängigkeiten können nachträglich geändert werden
	    \item DI ermöglicht flexible Konfiguration und Austausch von Komponenten
	    \item DI ist die Basis für Testbarkeit und saubere, modulare Architekturen
	\end{itemize}
\end{frame}

\begin{frame}{Vorteile von Dependency Injection}
	\begin{itemize}
	    \item Geringe Kopplung zwischen Komponenten
	    \item Leicht austauschbare Implementierungen
	    \item Fördert Testbarkeit (Mocking, Unit Tests)
	    \item Unterstützt flexible und zentrale Konfiguration (z.B. über Factories oder Frameworks)
	    \item Erleichtert Wartung und Erweiterung der Software
	\end{itemize}
\end{frame}

\begin{frame}{Mögliche Nachteile / Herausforderungen}
	\begin{itemize}
	    \item Bei vielen Abhängigkeiten kann der Konstruktor unübersichtlich werden
	    \item Setter Injection kann zu mutierbaren Objekten führen
	    \item Zirkuläre Abhängigkeiten können Probleme verursachen
	    \item Erhöhte Komplexität, wenn viele Fabriken und Konfigurationsklassen benötigt werden
	    \item Lernkurve für Entwickler, insbesondere ohne Framework-Unterstützung
	\end{itemize}
\end{frame}

\begin{frame}{Typische Anwendungsgebiete}
\begin{itemize}
    \item Frameworks wie Spring, Jakarta EE, Koin, Dagger
    \item Testbare Unit-Tests durch Mocking von Abhängigkeiten
    \item Flexible Konfiguration von Objekten bei Laufzeit oder Deployment
    \item Kombination mit anderen Mustern:
    \begin{itemize}
        \item \textbf{Strategy}: dynamischer Austausch von Verhalten
        \item \textbf{Factory}: zentrale Erzeugung und Konfiguration von Objekten
        \item \textbf{Singleton / Konfiguration}: zentrale Bereitstellung von Abhängigkeiten
    \end{itemize}
\end{itemize}
\end{frame}

\begin{frame}{Zusammenfassung}
	\begin{itemize}
	    \item Dependency Injection ist ein modernes Architekturpattern, kein klassisches GoF-Muster
	    \item Kernidee: \textbf{Abhängigkeiten werden von außen bereitgestellt}
	    \item Fördert lose Kopplung, Testbarkeit und Wiederverwendbarkeit
	    \item DI lässt sich gut mit Strategy, Factory und Singleton kombinieren
	    \item Constructor Injection ist bevorzugt, Setter Injection für optionale oder zirkuläre Abhängigkeiten
	\end{itemize}
\end{frame}
